\begin{proof}[Proof of \cref{obj:lem:randomized-kernel-barycenter}]
\begin{enumerate}
\item Fix a bounded kernel \(\kappa\) and an observation \(o\in\mathcal Z\). Its support condition gives
\[
\kappa(o,[-1,1])=1,
\qquad
-1\le h\le 1
\quad\text{for }\kappa(o,dh)\text{-almost every }h.
\]
The coordinate function is therefore integrable. Integrating these bounds against the probability measure \(\kappa(o,\cdot)\) yields
\[
-1=\int_{\mathbb R}(-1)\,\kappa(o,dh)
\le b_\kappa(o)
\le \int_{\mathbb R}1\,\kappa(o,dh)=1.
\]
Every map on the finite space \(\mathcal Z\) is measurable because its singletons are measurable. Thus \(b_\kappa\) is a measurable map into \([-1,1]\), determined by \(\kappa\). % lean: CausalSmith.Stat.MarRareqLogfrontier.kernel_barycenter_mem

\item Fix \(\lambda\in\Lambda\). Bounded support gives square integrability of the coordinate and its translate by \(\vartheta(\lambda)\), and
\[
\int_{\mathbb R}\bigl(h-\vartheta(\lambda)\bigr)\,\kappa(o,dh)
=b_\kappa(o)-\vartheta(\lambda).
\]
The second-moment formula for the variance of the translated coordinate, together with invariance of variance under subtraction of a constant, gives
\[
\begin{aligned}
\int_{\mathbb R}\bigl(h-b_\kappa(o)\bigr)^2\,\kappa(o,dh)
&=
\int_{\mathbb R}\bigl(h-\vartheta(\lambda)\bigr)^2\,\kappa(o,dh)\\
&\quad-
\left(
\int_{\mathbb R}\bigl(h-\vartheta(\lambda)\bigr)\,\kappa(o,dh)
\right)^2.
\end{aligned}
\]
Substituting the preceding mean identity and rearranging yields
\[
\int_{\mathbb R}\bigl(h-\vartheta(\lambda)\bigr)^2\,\kappa(o,dh)
=
\bigl(b_\kappa(o)-\vartheta(\lambda)\bigr)^2
+
\int_{\mathbb R}\bigl(h-b_\kappa(o)\bigr)^2\,\kappa(o,dh).
\]
% lean: CausalSmith.Stat.MarRareqLogfrontier.kernel_barycenter_variance_identity
The last integral is nonnegative, so
\[
\bigl(b_\kappa(o)-\vartheta(\lambda)\bigr)^2
\le
\int_{\mathbb R}\bigl(h-\vartheta(\lambda)\bigr)^2\,\kappa(o,dh).
\]
% lean: CausalSmith.Stat.MarRareqLogfrontier.kernel_barycenter_sq_le
Both sides are integrable as functions of \(o\), since \(\mathcal Z\) is finite and \(Q_\lambda\) is a probability measure. Integrating gives
\[
\int_{\mathcal Z}\bigl(b_\kappa(o)-\vartheta(\lambda)\bigr)^2\,Q_\lambda(do)
\le
\int_{\mathcal Z}\int_{\mathbb R}
\bigl(h-\vartheta(\lambda)\bigr)^2\,\kappa(o,dh)\,Q_\lambda(do).
\]
% lean: CausalSmith.Stat.MarRareqLogfrontier.kernel_barycenter_risk_le

\item Define the estimator classes and their real-valued risks by
\[
\begin{aligned}
\mathcal K_{\mathrm{bd}}
&:=
\left\{\kappa:\mathcal Z\leadsto\mathbb R:
\kappa\text{ is a parameter-independent Markov kernel and }
\kappa(o,[-1,1])=1\text{ for every }o\right\},\\
\mathcal B_{\mathrm{det}}
&:=\{b:\mathcal Z\to[-1,1]\},\\
\mathcal R_{\mathrm{bd}}(\kappa,\lambda)
&:=
\int_{\mathcal Z}\int_{\mathbb R}
\bigl(h-\vartheta(\lambda)\bigr)^2\,\kappa(o,dh)\,Q_\lambda(do),
\qquad (\kappa,\lambda)\in\mathcal K_{\mathrm{bd}}\times\Lambda,\\
\mathcal R_{\mathrm{det}}(b,\lambda)
&:=
\int_{\mathcal Z}\bigl(b(o)-\vartheta(\lambda)\bigr)^2\,Q_\lambda(do),
\qquad (b,\lambda)\in\mathcal B_{\mathrm{det}}\times\Lambda.
\end{aligned}
\]
All real suprema over an empty parameter set are taken to be zero.

For every bounded kernel, its support condition and the target bound imply
\[
-2\le h-\vartheta(\lambda)\le2,
\qquad
0\le\bigl(h-\vartheta(\lambda)\bigr)^2\le4
\quad\text{for }\kappa(o,dh)\text{-almost every }h.
\]
Integrating first against \(\kappa(o,\cdot)\) and then against \(Q_\lambda\) gives
\[
0\le
\int_{\mathbb R}\bigl(h-\vartheta(\lambda)\bigr)^2\,\kappa(o,dh)
\le4,
\qquad
0\le\mathcal R_{\mathrm{bd}}(\kappa,\lambda)\le4.
\]
% lean: CausalSmith.Stat.MarRareqLogfrontier.bounded_kernelRisk_bounds

For every \(b\in\mathcal B_{\mathrm{det}}\), define its deterministic kernel by
\[
\kappa_b(o,\cdot):=\delta_{b(o)},
\qquad o\in\mathcal Z.
\]
Measurability of \(b\) makes this a Markov kernel, and its interval-valued range gives
\[
\kappa_b(o,[-1,1])=1,
\qquad
\mathcal R_{\mathrm{bd}}(\kappa_b,\lambda)
=\mathcal R_{\mathrm{det}}(b,\lambda).
\]
% lean: CausalSmith.Stat.MarRareqLogfrontier.boundedMapKernel_risk
In particular, both estimator classes are nonempty: take
\[
b_{\mathrm{zero}}:\mathcal Z\to[-1,1],
\qquad b_{\mathrm{zero}}(o):=0,
\]
and its kernel \(\kappa_{b_{\mathrm{zero}}}\). The risk identity also gives
\(0\le\mathcal R_{\mathrm{det}}(b,\lambda)\le4\).
Consequently, all worst-case risks lie in \([0,4]\), including when
\(\Lambda\) is empty, and their infima are well defined.

The deterministic-kernel construction gives the first comparison:
\[
\inf_{\kappa\in\mathcal K_{\mathrm{bd}}}
\sup_{\lambda\in\Lambda}\mathcal R_{\mathrm{bd}}(\kappa,\lambda)
\le
\inf_{b\in\mathcal B_{\mathrm{det}}}
\sup_{\lambda\in\Lambda}\mathcal R_{\mathrm{det}}(b,\lambda).
\]
Conversely, each
\(\kappa\in\mathcal K_{\mathrm{bd}}\) has the barycenter map
\(b_\kappa\in\mathcal B_{\mathrm{det}}\), whose pointwise risk inequality gives
\[
\sup_{\lambda\in\Lambda}\mathcal R_{\mathrm{det}}(b_\kappa,\lambda)
\le
\sup_{\lambda\in\Lambda}\mathcal R_{\mathrm{bd}}(\kappa,\lambda).
\]
For a nonempty parameter set this follows from monotonicity of bounded
suprema; for an empty parameter set both sides equal zero. Taking infima gives
\[
\inf_{b\in\mathcal B_{\mathrm{det}}}
\sup_{\lambda\in\Lambda}\mathcal R_{\mathrm{det}}(b,\lambda)
\le
\inf_{\kappa\in\mathcal K_{\mathrm{bd}}}
\sup_{\lambda\in\Lambda}\mathcal R_{\mathrm{bd}}(\kappa,\lambda).
\]
The two comparisons establish the asserted equality. % lean: CausalSmith.Stat.MarRareqLogfrontier.bounded_kernel_minimax_eq_deterministic

\item Fix an arbitrary real-valued kernel \(\kappa\) and
\(\lambda\in\Lambda\). First suppose \(h\le1\). If \(-1\le h\), clipping fixes
\(h\), so the squared losses agree. If \(h\le-1\), clipping gives \(-1\), and
\[
\begin{aligned}
\bigl(h-\vartheta(\lambda)\bigr)^2
-\bigl(-1-\vartheta(\lambda)\bigr)^2
&=(-1-h)\bigl(1-h+2\vartheta(\lambda)\bigr)\\
&\ge0,
\end{aligned}
\]
because both factors are nonnegative. In the remaining case \(1\le h\),
clipping gives \(1\), and
\[
\begin{aligned}
\bigl(h-\vartheta(\lambda)\bigr)^2
-\bigl(1-\vartheta(\lambda)\bigr)^2
&=(h-1)\bigl(h+1-2\vartheta(\lambda)\bigr)\\
&\ge0.
\end{aligned}
\]
Thus, for every \(h\in\mathbb R\),
\[
\bigl(\operatorname{clip}_{[-1,1]}(h)-\vartheta(\lambda)\bigr)^2
\le
\bigl(h-\vartheta(\lambda)\bigr)^2.
\]
% lean: CausalSmith.Stat.MarRareqLogfrontier.realKernelClip_sq_le

Define the clipped kernel by the pushforward
\[
\kappa^{\mathrm{cl}}(o,\cdot)
:=
\kappa(o,\cdot)\circ\operatorname{clip}_{[-1,1]}^{-1},
\qquad o\in\mathcal Z,
\]
and define the extended risk for any real-valued Markov kernel by
\[
\mathcal R_{\mathrm{ext}}(\kappa,\lambda)
:=
\int_{\mathcal Z}\int_{\mathbb R}
\bigl(h-\vartheta(\lambda)\bigr)^2\,\kappa(o,dh)\,Q_\lambda(do)
\in[0,\infty],
\]
where both integrals are nonnegative extended integrals. Clipping is measurable.
The pushforward integration formula and monotonicity of these integrals give
\[
\begin{aligned}
\mathcal R_{\mathrm{ext}}(\kappa^{\mathrm{cl}},\lambda)
&=
\int_{\mathcal Z}\int_{\mathbb R}
\bigl(\operatorname{clip}_{[-1,1]}(h)-\vartheta(\lambda)\bigr)^2
\,\kappa(o,dh)\,Q_\lambda(do)\\
&\le \mathcal R_{\mathrm{ext}}(\kappa,\lambda).
\end{aligned}
\]
This comparison includes infinite original risks. % lean: CausalSmith.Stat.MarRareqLogfrontier.realKernelClip_risk_le

\item Define the full kernel class by
\[
\mathcal K_{\mathrm{real}}
:=
\{\kappa:\mathcal Z\leadsto\mathbb R:
\kappa\text{ is a parameter-independent Markov kernel}\}.
\]
For a bounded kernel, squared loss is nonnegative and integrable at every
observation. Its inner integral is integrable over the finite observation
space. Applying the identification of a nonnegative integrable function's real
and extended integrals at both levels gives
\[
\mathcal R_{\mathrm{ext}}(\kappa,\lambda)
=\mathcal R_{\mathrm{bd}}(\kappa,\lambda)
\quad\text{in }[0,\infty],
\qquad \kappa\in\mathcal K_{\mathrm{bd}}.
\]
% lean: CausalSmith.Stat.MarRareqLogfrontier.bounded_realKernelRisk_eq_ofReal

Inclusion of bounded kernels into the full class therefore gives
\[
\inf_{\kappa\in\mathcal K_{\mathrm{real}}}
\sup_{\lambda\in\Lambda}\mathcal R_{\mathrm{ext}}(\kappa,\lambda)
\le
\inf_{\kappa\in\mathcal K_{\mathrm{bd}}}
\sup_{\lambda\in\Lambda}\mathcal R_{\mathrm{bd}}(\kappa,\lambda),
\]
with the right side here evaluated in \([0,\infty]\).

For the reverse comparison, clipping takes values in \([-1,1]\), since
\[
-1\le \max\{-1,\min\{1,h\}\}\le1
\qquad(h\in\mathbb R).
\]
A measurable pushforward preserves probability, and
\[
\operatorname{clip}_{[-1,1]}^{-1}([-1,1])=\mathbb R,
\qquad
\kappa^{\mathrm{cl}}(o,[-1,1])=\kappa(o,\mathbb R)=1.
\]
Thus \(\kappa^{\mathrm{cl}}\in\mathcal K_{\mathrm{bd}}\) for every
\(\kappa\in\mathcal K_{\mathrm{real}}\). The preceding risk comparison yields
\[
\inf_{\kappa\in\mathcal K_{\mathrm{bd}}}
\sup_{\lambda\in\Lambda}\mathcal R_{\mathrm{bd}}(\kappa,\lambda)
\le
\inf_{\kappa\in\mathcal K_{\mathrm{real}}}
\sup_{\lambda\in\Lambda}\mathcal R_{\mathrm{ext}}(\kappa,\lambda).
\]
Hence the two extended minimax values agree.

For the conversion to a real value, use
\[
\operatorname{rv}:[0,\infty]\to\mathbb R,
\qquad
\operatorname{rv}|_{[0,\infty)}=\operatorname{id},
\qquad
\operatorname{rv}(\infty)=0.
\]
The bounds \(0\le\mathcal R_{\mathrm{bd}}\le4\) place every bounded-kernel
supremum, and their infimum, in \([0,4]\). On this interval the real and
extended orders agree, so conversion commutes with these suprema and the
infimum. For an empty parameter set both kinds of supremum equal zero.
It follows that
\[
\operatorname{rv}\!\left(
\inf_{\kappa\in\mathcal K_{\mathrm{real}}}
\sup_{\lambda\in\Lambda}\mathcal R_{\mathrm{ext}}(\kappa,\lambda)
\right)
=
\inf_{\kappa\in\mathcal K_{\mathrm{bd}}}
\sup_{\lambda\in\Lambda}\mathcal R_{\mathrm{bd}}(\kappa,\lambda),
\]
where the right side is now evaluated in \(\mathbb R\). % lean: CausalSmith.Stat.MarRareqLogfrontier.real_kernel_minimax_eq_bounded

\item Define the seed law, viewed as a measure on \([0,1]\), by
\[
\nu_{\mathrm{seed}}
:=\left.\operatorname{Leb}_{\mathbb R}\right|_{[0,1]}.
\]
Its total mass is
\[
\nu_{\mathrm{seed}}([0,1])
=\operatorname{Leb}_{\mathbb R}([0,1])
=1-0=1,
\]
so it is a probability law. % lean: CausalSmith.Stat.MarRareqLogfrontier.uniformSeedLaw_probability

\item Fix \(\kappa\in\mathcal K_{\mathrm{bd}}\). The measurable randomization
theorem for Markov kernels with a standard Borel output space, applied to the
real output space, supplies a jointly measurable map satisfying
\[
\widetilde f:\mathcal Z\times[0,1]\to\mathbb R,
\qquad
\kappa(o,\cdot)
=
\nu_{\mathrm{seed}}\circ\widetilde f(o,\cdot)^{-1}
\quad(o\in\mathcal Z).
\]
Define the interval-valued realization by
\[
f:\mathcal Z\times[0,1]\to[-1,1],
\qquad
f(o,v):=\operatorname{clip}_{[-1,1]}\bigl(\widetilde f(o,v)\bigr).
\]
It is jointly measurable by measurability of clipping. For
\(h\in[-1,1]\), clipping fixes \(h\); the bounded support condition consequently
gives
\[
\operatorname{clip}_{[-1,1]}(h)=h
\quad\text{for }\kappa(o,dh)\text{-almost every }h,
\qquad
\kappa(o,\cdot)\circ\operatorname{clip}_{[-1,1]}^{-1}
=\kappa(o,\cdot).
\]
% lean: CausalSmith.Stat.MarRareqLogfrontier.boundedKernel_map_clip
Composition of measurable pushforwards now yields
\[
\begin{aligned}
\nu_{\mathrm{seed}}\circ f(o,\cdot)^{-1}
&=
\bigl(\nu_{\mathrm{seed}}\circ\widetilde f(o,\cdot)^{-1}\bigr)
\circ\operatorname{clip}_{[-1,1]}^{-1}\\
&=\kappa(o,\cdot)\circ\operatorname{clip}_{[-1,1]}^{-1}
=\kappa(o,\cdot).
\end{aligned}
\]
Thus \(f(o,U)\), for \(U\sim\nu_{\mathrm{seed}}\), has the required distribution
at every observation. % lean: CausalSmith.Stat.MarRareqLogfrontier.boundedKernel_uniformSeed_representation

\item Conversely, let \(f:\mathcal Z\times[0,1]\to[-1,1]\) be jointly
measurable, and define
\[
\kappa_f(o,\cdot)
:=
\nu_{\mathrm{seed}}\circ f(o,\cdot)^{-1},
\qquad o\in\mathcal Z.
\]
Every section \(f(o,\cdot)\) is measurable, so its pushforward is a probability
measure. Finiteness of \(\mathcal Z\) and measurability of its singletons give
the kernel measurability condition. Moreover,
\[
f(o,\cdot)^{-1}([-1,1])=[0,1],
\qquad
\kappa_f(o,[-1,1])
=\nu_{\mathrm{seed}}([0,1])=1.
\]
Hence \(\kappa_f\) is a bounded Markov kernel with exactly the displayed
conditional distributions. Both constructions use the same law
\(\nu_{\mathrm{seed}}\), independently of the parameter. For each
\(\lambda\), the auxiliary seed can be sampled independently of the
observation under the product law
\[
Q_\lambda\otimes\nu_{\mathrm{seed}}.
\]
% lean: CausalSmith.Stat.MarRareqLogfrontier.boundedKernel_of_uniformSeed
\end{enumerate}
\end{proof}
